\section*{Вступ}

Дисципліна «Програмування» є базовою складовою підготовки студентів спеціальності 
Е7 «Математика» та формує фундаментальні знання і навички, необхідні для подальшого 
застосування сучасних інформаційних технологій у наукових дослідженнях та професійній діяльності. 
Вона забезпечує розвиток алгоритмічного мислення, вміння будувати та аналізувати алгоритми, 
а також практичні навички роботи з мовами програмування.

Особливе місце у курсі займає мова Python, яка завдяки простоті синтаксису, універсальності 
та широкому спектру бібліотек стала одним із ключових інструментів у сфері аналізу даних, 
машинного навчання та Data Science. Для студентів-математиків це відкриває можливості 
поєднання теоретичних знань з практичними застосуваннями у статистиці, моделюванні, 
обробці великих масивів даних та побудові прогнозних моделей.

Основними завданнями курсу є:
\begin{itemize}
	\item ознайомлення студентів з базовими алгоритмічними конструкціями структурного програмування;
	\item формування навичок написання програм мовою Python;
	\item розвиток умінь застосовувати програмування для розв’язання математичних задач;
	\item підготовка до використання Python у подальших дисциплінах, пов’язаних з аналізом даних та Data Science.
\end{itemize}
Таким чином, курс «Програмування» є важливим етапом у професійній підготовці студентів 
спеціальності Е7 «Математика», оскільки створює основу для подальшого опанування сучасних 
методів аналізу даних, статистичних моделей та технологій штучного інтелекту.